\newpage
\section{CUDA}
\label{sec:CUDA}
I implemented a CUDA version of the Mandelbrot computation and compared it with a CPU baseline, based on the results obtained with my PC.
\noindent
I defined the parameters as follows: $RESOLUTION = 5000$ and $ITERATIONS = 100$, resulting in an image of $15{,}000 \times 10{,}000$ pixels (approximately 150 million points).

\subsection{Code}

\paragraph{Memory allocation:} this step prepares the memory needed for the Mandelbrot output.
\begin{lstlisting}
    const int N = WIDTH * HEIGHT;
    int *image_host = new int[N];
    int *image_dev;

    cudaMalloc(&image_dev, N * sizeof(int));
\end{lstlisting}

\paragraph{Kernel configuration and execution:} the code defines the CUDA execution parameters: a block of 256 threads and a 1D grid large enough to cover all $N$ pixels.
\begin{lstlisting}
    dim3 block(256);
    dim3 grid((N + block.x - 1) / block.x);

    mandelbrot_kernel_1D<<<grid, block>>>(
        image_dev, WIDTH, HEIGHT, STEP, ITERATIONS, MIN_X, MIN_Y
    );
    
    cudaDeviceSynchronize();
\end{lstlisting}

\paragraph{Kernel Mandelbrot 1D:} the \texttt{mandelbrot\_kernel\_1D()} is a GPU kernel that parallelizes the Mandelbrot computation by assigning one thread per pixel.

\begin{lstlisting}[basicstyle=\ttfamily\scriptsize]
__global__ void mandelbrot_kernel_1D(int *image, int width, int height,
                                  double step, int iterations,
                                  double min_x, double min_y) {
    int pos = blockIdx.x * blockDim.x + threadIdx.x;
    int size = width * height;
    if (pos >= size) return;

    int row = pos / width;
    int col = pos % width;
    double cRe = col * step + min_x;
    double cIm = row * step + min_y;
    double zRe = 0, zIm = 0;

    for (int i = 1; i <= iterations; i++) {
        double oldRe = zRe;
        zRe = (zRe*zRe - zIm*zIm) + cRe;
        zIm = 2 * oldRe * zIm + cIm;

        if (zRe*zRe + zIm*zIm >= 4.0) {
            image[pos] = i;
            return;
        }
    }

    image[pos] = 0;
}
\end{lstlisting}

\paragraph{Copy-back and cleanup:} After the GPU finishes, the results are copied from device to host using:
\begin{lstlisting}
    cudaMemcpy(image_host, image_dev, N * sizeof(int), cudaMemcpyDeviceToHost);

    cudaFree(image_dev);
\end{lstlisting}

\subsection{Compilation}

The compilation commands used for the different architectures were the following:

\paragraph{Personal PC:}
\begin{verbatim}
    CPU: g++ -O2 -o build/opt1 src/opt1.cpp
    GPU: nvcc -O2 -arch=sm_50 src/*.cu -o build/*
\end{verbatim}
\paragraph{Google Colab:}
\begin{verbatim}
    !nvcc -O2 -arch=sm_75 -o "$ROOT/build/*" "$ROOT/src/*.cu"
\end{verbatim}

\subsection{Results}
Both 1D and 2D CUDA kernels were tested, but they produced identical results, so I considered only the 1D version.\\

\noindent
Additionally, a version of the kernel using \textbf{float precision} (code in: \texttt{src/mandelbrot\_float.cu}) was tested, providing a significant speedup on GPU. The float kernel uses \textit{float} instead of \textit{double} for all real-number calculations, reducing memory usage and improving throughput on the GPU.

\begin{table}[h!]
\centering
\begin{tabular}{lccc}
\hline
\textbf{Kernel} & \textbf{CPU} & \textbf{GPU 950M} & \textbf{GPU T4} \\
\hline
1D        & 24668 ms & 1768 ms (13x) & 506 ms (48x) \\
Float     & -        & 194 ms (127x) & 27.8 ms (887x) \\
\hline
\end{tabular}
\end{table}